\DocumentMetadata{
  pdfversion=2.0,
  pdfstandard=ua-2,
  lang=en-US,
  tagging=on,
  tagging-setup={math/setup=mathml-SE,tabsorder=structure}
}
\documentclass[11pt,letterpaper]{article}
\usepackage[margin=0.68in,includefoot]{geometry}
\usepackage{newcomputermodern}
\usepackage{amsmath,amscd,mathtools}
\usepackage{tikz,adjustbox}
\providecommand{\square}{\mdlgwhtsquare}
\usepackage[hidelinks]{hyperref}
\hypersetup{
  pdftitle={Numerical Analysis PhD Exam, January 2002},
  pdfauthor={University of Florida Department of Mathematics},
  pdfsubject={Graduate mathematics examination},
  pdfdisplaydoctitle=true,
  bookmarksopen=true
}
\setcounter{secnumdepth}{0}
\setlength{\parindent}{0pt}
\setlength{\parskip}{0.28em}
\setlength{\emergencystretch}{3em}
\setlength{\leftmargini}{2.15em}
\setlength{\leftmarginii}{2.65em}
\setlength{\leftmarginiii}{3em}
\pagestyle{empty}
\raggedbottom
\allowdisplaybreaks
\definecolor{ProblemConcern}{HTML}{7A1F1F}
\newenvironment{problemconcern}{%
  \par\tagstructbegin{tag=Aside}\begingroup\color{ProblemConcern}
}{%
  \par\endgroup\tagstructend\nopagebreak[4]\vspace{0.12em}
}
\NewTaggingSocketPlug{block/list/label}{examproblem}{%
  \tagstructbegin{tag=Lbl}\tagstructend%
  \tagstructbegin{tag=\LBody}%
  \tagstructbegin{tag=\ProblemHeadingTag,title-o={\ProblemHeadingText}}%
  \tagmcbegin{tag=\ProblemHeadingTag,actualtext={\ProblemHeadingText}}%
  #2%
  \tagmcend%
  \tagstructend%
}
\newenvironment{examproblems}[2]{%
  \def\ProblemHeadingTag{#2}%
  \AssignTaggingSocketPlug{block/list/label}{examproblem}%
  \begin{enumerate}%
  \setcounter{enumi}{#1}%
  \renewcommand{\labelenumi}{\textbf{\theenumi.}}%
  \setlength{\topsep}{0.35em}%
  \setlength{\partopsep}{0pt}%
  \setlength{\itemsep}{0.48em}%
  \setlength{\parsep}{0.18em}%
}{\end{enumerate}}
\newenvironment{examsubparts}{%
  \AssignTaggingSocketPlug{block/list/label}{default}%
  \begin{enumerate}%
  \setlength{\topsep}{0.18em}%
  \setlength{\partopsep}{0pt}%
  \setlength{\itemsep}{0.16em}%
  \setlength{\parsep}{0.1em}%
}{\end{enumerate}}
\newenvironment{examinstructions}{%
  \par\begingroup\itshape
}{%
  \par\endgroup\vspace{0.18em}
}
\makeatletter
\renewcommand\subsection{\@startsection{subsection}{2}{0pt}%
  {-1.25ex plus -.25ex minus -.1ex}%
  {0.55ex plus .1ex}%
  {\normalfont\large\bfseries}}
\renewcommand{\@maketitle}{%
  \newpage\null\vskip -1em%
  {\footnotesize\bfseries
    \href{https://www.ufl.edu/}{University of Florida}, \href{https://math.ufl.edu/}{Department of Mathematics}\par}%
  \vskip -0.5em%
  \tagstructbegin{tag=H1,title={Numerical Analysis PhD Exam, January 2002}}%
  \tagpdfparaOff%
  \tagmcbegin{tag=H1}%
    {\LARGE\bfseries\href{https://gma.math.ufl.edu/exams/numerical-analysis-phd/}{Numerical Analysis PhD Exam}, January 2002\par}%
  \tagmcend%
  \tagpdfparaOn%
  \tagstructend%
  \vskip 0.25em%
  \tagmcbegin{artifact}%
    \hrule height 1pt%
  \tagmcend%
  \vskip 1em%
}
\makeatother
\title{Numerical Analysis PhD Exam, January 2002}
\author{}
\date{}
\begin{document}
\maketitle
\thispagestyle{empty}
\begin{examinstructions}
Do any 8 of the following 10 problems. Problems 1–5 are Part 1: Numerical Linear Algebra. Problems 6–10 are Part II: Numerical Analysis.
\end{examinstructions}
\begin{examproblems}{0}{H2}
\def\ProblemHeadingText{Problem 1.}
\item
This problem has three parts.
\begin{examsubparts}
\item[\textnormal{(a)}]
Prove: If \(A\in\mathbb R^{n\times n}\) has a set of~\(n\) linearly independent eigenvectors, then~\(A\) can be diagonalized.
\item[\textnormal{(b)}]
Prove: If \(A\in\mathbb R^{n\times n}\) is symmetric, then~\(A\) has an orthogonal set of eigenvectors.
\item[\textnormal{(c)}]
Diagonalize the matrix
\[
A=\begin{pmatrix}66&12\\12&59\end{pmatrix}.
\]
\end{examsubparts}
\def\ProblemHeadingText{Problem 2.}
\item
This problem has three parts.
\begin{examsubparts}
\item[\textnormal{(a)}]
Is the matrix
\[
A=\begin{pmatrix}5&1\\0&3\end{pmatrix}
\]
 diagonalizable? Why?
\item[\textnormal{(b)}]
Can the above matrix~\(A\) be diagonalized by an orthogonal matrix~\(S\)?
\item[\textnormal{(c)}]
State a theorem or condition that ensures a matrix can be diagonalized by an orthogonal matrix.
\end{examsubparts}
\def\ProblemHeadingText{Problem 3.}
\item
This problem has three parts.
\begin{examsubparts}
\item[\textnormal{(a)}]
Give a careful statement of the singular value decomposition.
\item[\textnormal{(b)}]
Give a careful proof of the singular value decomposition.
\item[\textnormal{(c)}]
Compute the SVD for the matrix
\[
A=\begin{pmatrix}-6&17\\18&-1\end{pmatrix}.
\]
\end{examsubparts}
\def\ProblemHeadingText{Problem 4.}
\item
\begin{problemconcern}
\textbf{Warning: Suspected error.} The source clearly asks for the projection of~\(A\) onto its own column space, but the intended object is mathematically ambiguous; it may have intended the orthogonal projection matrix onto the column space. The source wording has been retained.
\end{problemconcern}
Let
\[
A=\begin{pmatrix}1&2\\0&1\\1&0\end{pmatrix}.
\]
\begin{examsubparts}
\item[\textnormal{(a)}]
Determine the orthogonal projection of~\(A\) onto the column space of~\(A\).
\item[\textnormal{(b)}]
Compute the \(QR\) factorization of~\(A\).
\end{examsubparts}
\def\ProblemHeadingText{Problem 5.}
\item
\begin{problemconcern}
\textbf{Warning: Suspected error.} The displayed dimensions make the product \(BA\) undefined in general. The source reading has been retained because it is unclear whether the dimensions or the product order should be corrected.
\end{problemconcern}
This problem has five parts.
\begin{examsubparts}
\item[\textnormal{(a)}]
Define the operator 2-norm for any matrix~\(A\) in \(\mathbb R^{m\times n}\).
\item[\textnormal{(b)}]
Prove: If \(A\in\mathbb R^{m\times n}\) and \(B\in\mathbb R^{n\times q}\), then \(\lVert BA\rVert_2\leq\lVert B\rVert_2\cdot\lVert A\rVert_2\).
\item[\textnormal{(c)}]
Prove: If~\(A\) is any matrix in \(\mathbb R^{n\times n}\) and \(\lVert A\rVert_2<1\), then \(\lim_{n\to\infty}A^n=0\).
\item[\textnormal{(d)}]
Prove: If~\(A\) is a matrix in \(\mathbb R^{n\times n}\), then define \(e^A\) as a power series expansion. Show this series converges.
\item[\textnormal{(e)}]
Prove: If \(A\in\mathbb R^{n\times n}\) is invertible, \(b\) and \(\Delta b\) are vectors in \(\mathbb R^n\), \(x\) and \(\Delta x\) are vectors in \(\mathbb R^n\) and are solutions to \(Ax=b\) and \(A\Delta x=\Delta b\), respectively, then
\[
\frac{\lVert\Delta x\rVert_2}{\lVert x\rVert_2}\leq\kappa_2(A)\frac{\lVert\Delta b\rVert_2}{\lVert b\rVert_2}.
\]
\end{examsubparts}
\def\ProblemHeadingText{Problem 6.}
\item
This problem has three parts.
\begin{examsubparts}
\item[\textnormal{(a)}]
State Newton’s method (the algorithm) for solving \(f(x)=0\) where \(f:\mathbb R\to\mathbb R\).
\item[\textnormal{(b)}]
Assuming~\(f\) is smooth and we have an initial guess sufficiently close to the root, state sufficient condition(s) for the method to converge quadratically.
\item[\textnormal{(c)}]
Prove: Let \(T(x):[a,b]\to[a,b]\) be a function such that \(T'''(x)\) is continuous for all \(x\in[a,b]\), \(T(p)=p\), \(T'(p)=T''(p)=0\). If \(x_0\in[a,b]\) is a point with the property that the sequence defined recursively by \(x_{k+1}=T(x_k)\) converges to~\(p\), then the sequence \(\{x_k\}_{k=1}^{\infty}\) converges cubically to~\(p\).
\end{examsubparts}
\def\ProblemHeadingText{Problem 7.}
\item
This problem has four parts.
\begin{examsubparts}
\item[\textnormal{(a)}]
Give a careful statement of the minimization property for clamped cubic splines.
\item[\textnormal{(b)}]
Give a careful statement of the convergence theorem for clamped cubic splines.
\item[\textnormal{(c)}]
Give a careful statement of the convergence theorem for the second derivatives for clamped cubic splines.
\item[\textnormal{(d)}]
Outline a proof of the convergence theorem for clamped cubic splines.
\end{examsubparts}
\def\ProblemHeadingText{Problem 8.}
\item
This problem has two parts.
\begin{examsubparts}
\item[\textnormal{(a)}]
Given a set of data points \((x_0,y_0),(x_1,y_1),\ldots,(x_n,y_n)\) in the plane, describe how to find the conic section of the form \(Ax^2+Bxy+Cy^2+Dx+Ey+1=0\) which best fits the data in the sense of least squares.
\item[\textnormal{(b)}]
How do you determine whether or not the conic section is an ellipse, a hyperbola, or a parabola?
\end{examsubparts}
\def\ProblemHeadingText{Problem 9.}
\item
Triple Recursion Formula. If \(\{\phi_n(x)\}\) is an orthogonal family of polynomials on \([a,b]\), with respect to the weight function \(w(x)\geq0\), and \(n\geq1\), then show that \(\phi_{n+1}(x)=(a_nx+b_n)\phi_n(x)-c_n\phi_{n-1}(x)\), for some constants \(a_n,b_n,\) and \(c_n\). (Hint: Let \(g(x)=\phi_{n+1}(x)-a_nx\phi_n(x)\), where \(a_n\) is a constant chosen so that \(g(x)\) has degree~\(n\).)
\def\ProblemHeadingText{Problem 10.}
\item
This problem has four parts.
\begin{examsubparts}
\item[\textnormal{(a)}]
Give a careful statement of the contraction mapping theorem.
\item[\textnormal{(b)}]
Prove the contraction mapping theorem.
\item[\textnormal{(c)}]
Determine whether or not the function \(T(x)=\frac14\sin(3x+2)+12\) is a contraction.
\item[\textnormal{(d)}]
Discuss the convergence rate ensured by the contraction mapping theorem. (ie Is it linear, quadratic, or cubic?)
\end{examsubparts}
\end{examproblems}
\end{document}
